% --------------------------------------
% Master's Thesis Title Page
% LaTeX Template
% Version 1.0 (23/05/14)
% Thanks to Magnus Marthinsen, this thesis and template is made available for Master studentes at HVL Joint SE program (01.03.2021)
%---------------------------------------

%----------------------------------------------------------------------------------------
%	PACKAGES AND OTHER DOCUMENT CONFIGURATIONS
%----------------------------------------------------------------------------------------
\documentclass[a4paper]{report}
\usepackage{graphicx} % Required for box manipulation
\usepackage{helvet}
\usepackage{subfig}
\usepackage[utf8]{inputenc}
%\usepackage{natbib}
\usepackage[USenglish]{babel}
\usepackage[useregional]{datetime2}
\usepackage{pgfgantt}
\usepackage{listings}
\usepackage{wrapfig}
\usepackage{setspace}
\usepackage{parskip} % Used to create spaces between paragraphs
\usepackage{dirtytalk} % quotes by talk
\usepackage[hidelinks]{hyperref}
% \usepackage[acronym, toc]{glossaries}  % Used to add a wordlist/glossaries
\usepackage{mathtools}
\usepackage{amsfonts}
\usepackage{color, colortbl}
\usepackage{booktabs}
\usepackage{float}
\usepackage{csquotes}

%BIB by Adrian
\usepackage[backend=biber,style=numeric, urldate=long]{biblatex}
% See the references.bib file. Most Bibtex bibliographies in Computer Science can be found from dblp.org 
\addbibresource{references.bib}



% Glossary/wordlist
% \makeglossaries
% \input{glossaries.tex}

\begin{document}

%
% COLORS USED THROUGH THE REPORT
%
\definecolor{kb_red}{RGB}{96,2,4}
\definecolor{light_gray}{RGB}{160,160,160}
\definecolor{med_gray}{RGB}{96,94,94}
\definecolor{black}{RGB}{0,0,0}
\definecolor{white}{RGB}{155,155,155}
\definecolor{light_green}{RGB}{208,240,192}
\definecolor{light_red}{RGB}{255,204,203}

% CODE STYLE
\definecolor{javared}{rgb}{0.6,0,0} % for strings
\definecolor{javagreen}{rgb}{0.25,0.5,0.35} % comments
\definecolor{javapurple}{rgb}{0.5,0,0.35} % keywords
\definecolor{javadocblue}{rgb}{0.25,0.35,0.75} % javadoc

\makeatletter
\lst@Key{matchrangestart}{f}{\lstKV@SetIf{#1}\lst@ifmatchrangestart}
\def\lst@SkipToFirst{%
    \lst@ifmatchrangestart\c@lstnumber=\numexpr-1+\lst@firstline\fi
    \ifnum \lst@lineno<\lst@firstline
        \def\lst@next{\lst@BeginDropInput\lst@Pmode
        \lst@Let{13}\lst@MSkipToFirst
        \lst@Let{10}\lst@MSkipToFirst}%
        \expandafter\lst@next
    \else
        \expandafter\lst@BOLGobble
    \fi}
\makeatother

\lstset{language=Java,
basicstyle=\footnotesize,
keywordstyle=\color{javapurple}\bfseries,
stringstyle=\color{javared},
commentstyle=\color{javagreen},
morecomment=[s][\color{javadocblue}]{/**}{*/},
numbers=left,
captionpos=b,
frame=single,
breakatwhitespace=false,
breaklines=true,
numberstyle=\tiny\color{black},
stepnumber=1,
numbersep=10pt,
tabsize=2,
showspaces=false,
showstringspaces=false,
matchrangestart=t}

%Title page
%----------------------------------------------------------------------------------------
%	TITLE PAGE
%----------------------------------------------------------------------------------------

\newcommand*{\titlePage}{\begingroup % Create the command for including the title page in the document
\fontfamily{phv}\selectfont
\centering % Center all text
\DTMlangsetup[en-US]{showdayofmonth=false}

%----------------------------------------------------------------------------------------
%	TITLE SECTION
%----------------------------------------------------------------------------------------

\vspace{200pt}
{\Huge VR viewer} \\ % Title
\vspace{5pt}

{\Large \textsl{}} % Subtitle or further description
\vspace{50pt}

%----------------------------------------------------------------------------------------
%	AUTHOR SECTION
%----------------------------------------------------------------------------------------

\LARGE{\textbf{Tobias Eilertsen}}\\ % Author name

\vfill % Whitespace between the author name and the school

%----------------------------------------------------------------------------------------
%	DESCRIPTION AND DATE SECTION
%----------------------------------------------------------------------------------------


{\Large \textbf{Master's thesis in Software Engineering at} \\
\vspace{10pt}
Department of Computer science, Electrical engineering and Mathematical sciences, \\
Western Norway University of Applied Sciences \\
\vspace{10pt}
Department of  Informatics, \\
University of Bergen \\}
\vspace{10pt}
{\large \today} % Month and year published

%----------------------------------------------------------------------------------------
%	LOGO SECTION
%----------------------------------------------------------------------------------------

\vfill % Whitespace between the school and the publisher logo


\begin{figure}[h!]
	\captionsetup[subfigure]{labelformat=empty}
	\subfloat[][]{\includegraphics[height=70pt]{images/logos/hvl_logo_engelsk.pdf}}
	\hfill
	\subfloat[][]{\includegraphics[height=70pt]{images/logos/uib-logo.eps}}
\end{figure}


\endgroup}

\titlePage
\pagebreak

\section*{Abstract}
% \acrlong{MDSE} is a ...

\section*{Acknowledgements}
First and foremost, I would like to thank ...
\pagebreak
\tableofcontents
\listoffigures
\listoftables

% \printglossary[nonumberlist]
% \printglossary[type=\acronymtype, nonumberlist]


\chapter{Introduction}
% \acrfull{SE} is an engineering discipline that focuses on the development of high-quality software systems \cite{TsuiFrankF2011Eose}. ...

\section{Context and Approach}
\section{Problem Description}
\section{Methodology}
\section{Contribution}
\section{Outline}

\chapter{Background} \label{Background}
In this chapter, we will present some of the knowledge that our research is built upon. This theory is important to know in order understand the following chapters.
...

$<$ \textit{The following sections and subsections are just examples of how to structure the background} $>$
\section{Model Driven Software Engineering}
\subsection{Modeling languages}\label{ModelingLanguages}

\section{Machine Learning} \label{Machine Learning}
\subsection{Supervised Learning}
\subsection{Unsupervised Learning}
\subsection{Reinforcement Learning}
\subsubsection{Q-learnng} \label{Q-learning}



\chapter{Design and Implementation} \label{Design and Implementation}
In this chapter the implementation of the algorithm will be explained. ...
\section{Demonstration} \label{demonstration}
\section{Development method}
\section{Code structure} \label{CodeStructure}


\chapter{Use cases} \label{Use Cases}
\chapter{Analysis and Assessment}\label{Analysis and Assessment}
\chapter{Discussion} \label{Discussion}
\chapter{Related Work} \label{Related Work}
\chapter{Conclusion} \label{Conclusion}
\chapter{Further Work} \label{Further Work}

\appendix
\chapter{Source code} \label{SourceCode}

The source code for the VR application is available at this URL: \url{https://github.com/tobias2912/VR-DICOM-viewer}.



%\bibliographystyle{splncs04}
%\bibliography{references}
\printbibliography

\end{document}i
